% ch1_b §1.3 (书页 9–14 / p023–p028)
%--A-TAIL--
%（以下为 §1.2 的收尾段落，印于书页 9 顶部，接式 (1.22) 之后）
证明只需把整个级数乘以 $\exp(-i\boldsymbol{g}\cdot\mathbf{r})$，然后积分。若 $\boldsymbol{g}$ 具有 (1.20) 的形式，则 $\exp(i\boldsymbol{g}\cdot\mathbf{r})$ 在一个晶胞上的积分显然为零，除非 $\boldsymbol{g}=0$。

\section{倒格子的性质}

由 (1.20) 定义的矢量，即
\begin{equation*}
\boldsymbol{g} = n_1\cdot 2\pi\mathbf{b}_1 + n_2\cdot 2\pi\mathbf{b}_2 + n_3\cdot 2\pi\mathbf{b}_3,
\tag{1.23}
\end{equation*}
生成一个点阵，其晶胞由矢量 $2\pi\mathbf{b}_1$、$2\pi\mathbf{b}_2$、$2\pi\mathbf{b}_3$ 张成。这就是我们原来正格子（direct lattice）的倒格子（reciprocal lattice）。它是一个不变的几何客体，其性质是固体理论的基础。它的几个初等几何性质很容易推导。

(i) \emph{倒格子的每一个矢量都垂直于正格子的一族点阵面（lattice planes）。}

任取一个倒格矢 $\boldsymbol{g}$ 和一个格矢 $\boldsymbol{l}$，考虑关系式 (1.21)：
\begin{align*}
\boldsymbol{g}\cdot\boldsymbol{l} &= 2\pi(n_1 l_1 + n_2 l_2 + n_3 l_3) \\
&= 2\pi N,
\tag{1.24}
\end{align*}
其中 $N$ 是整数。这告诉我们，矢量 $\boldsymbol{l}$ 在 $\boldsymbol{g}$ 方向上的投影长度为
\begin{equation*}
d = \frac{2\pi N}{|\boldsymbol{g}|}.
\tag{1.25}
\end{equation*}
但是，具有这一性质的格点有无穷多个。例如，设 $\boldsymbol{l}'$ 是由整数
\begin{equation*}
l'_1 = l_1 - mn_3, \qquad l'_2 = l_2 - mn_3, \qquad l'_3 = l_3 + m(n_1+n_2),
\tag{1.26}
\end{equation*}
所代表的格点，其中 $m$ 是整数（译注：原书第三式误印作 $l_1+m(n_1+n_2)$，按上下文应为 $l_3$）。显然
\begin{equation*}
\boldsymbol{g}\cdot\boldsymbol{l}' = \boldsymbol{g}\cdot\boldsymbol{l} = 2\pi N.
\tag{1.27}
\end{equation*}
于是 $\boldsymbol{l}'$ 在 $\boldsymbol{g}$ 上的投影与 $\boldsymbol{l}$ 相同，因而必定也位于与 $\boldsymbol{g}$ 垂直、距原点为 $d$ 的平面上。这样，若这个平面上有一个格点，其上就有无穷多个格点；我们便构造出了一个\emph{点阵面}。正格子与倒格子之间的这种关系，当然是三维几何中人所共知的点与平面对偶性的一个特例。

(ii) \emph{若 $\boldsymbol{g}$ 的分量没有公因子，则 $|\boldsymbol{g}|$ 与垂直于 $\boldsymbol{g}$ 的点阵面的间距成反比。}

这由式 (1.25) 可得（见图 8）。若 $(n_1, n_2, n_3)$ 没有公因子，则总能找到一个格矢 $\boldsymbol{l}''$，其分量满足
\begin{equation*}
\boldsymbol{g}\cdot\boldsymbol{l}'' = 2\pi(N+1).
\tag{1.28}
\end{equation*}
这意味着，包含 $\boldsymbol{l}''$ 的点阵面位于距原点
\begin{equation*}
d'' = \frac{2\pi(N+1)}{|\boldsymbol{g}|}
\tag{1.29}
\end{equation*}
处——亦即它与包含 $\boldsymbol{l}$ 的点阵面相距 $2\pi/|\boldsymbol{g}|$。

%FIG{8}: 图 8

从这两个初等几何结果可以看出，表征点阵诸平面最简单的方式就是用它们的法线，把它们表示为倒格矢。正格子中最显著的平面，是格点最密集的那些平面。由于格点在空间中的密度是常数，最显著的平面必定是相隔最远的那些平面，即具有最小倒格矢的那些平面。

用相应的倒格矢来标示点阵面，等价于经典结晶学中所用的米勒指数（Miller indices）。设有一个点阵面，法线为 $\boldsymbol{g}$，其上所有的点 $\boldsymbol{l}$ 都满足 (1.24)。取 $l_2 = l_3 = 0$ 的点，则得 $l_1 = N/n_1$，于是这个平面沿 $\mathbf{a}_1$ 轴的截距长度为
\begin{equation*}
d_1 = \left(\frac{N}{n_1}\right) a_1.
\tag{1.30}
\end{equation*}
类似地，这个平面与 $\mathbf{a}_2$ 轴相交于距原点
\begin{equation*}
d_2 = \left(\frac{N}{n_2}\right) a_2
\tag{1.31}
\end{equation*}
处。该平面沿各轴的截距，以相应基矢的长度为单位来量度，与整数 $n_1$、$n_2$、$n_3$ 成反比。除去公因子之后，这些整数就是这个点阵面的米勒指数，写成 $(n_1, n_2, n_3)$ 的形式。

显然，格点密集的面——即米勒指数小的面——最容易在天然晶体中出现，无论是在生长时还是解理之后。研究这类晶面的几何曾是经典结晶学的精髓，而其
%FIG{9}: 图 9　具有大米勒指数的晶面，其间距比主对称面的间距更小。
关键性的发现是：可以找到单独一组矢量 $\mathbf{a}_1$、$\mathbf{a}_2$、$\mathbf{a}_3$，使得宏观晶体上所有观察到的晶面都能用小的米勒指数来表示。

应当提到两种记号约定。符号 $(1\bar{1}\bar{1})$ 代表 $(1,-1,-1)$ 这组平面；为简洁起见，把负号写在数字上方。带花括号的符号 $\{110\}$ 代表在对称性上与 $(110)$ 等价的各组不同的平面。因此，对立方结构而言，它可以包括 $(101)$、$(011)$、$(\bar{1}\bar{1}0)$、$(1\bar{1}0)$，等等。

(iii) \emph{倒格子晶胞的体积与正格子晶胞的体积成反比。}

这用初等矢量分析即可得出。倒格子晶胞由 $2\pi\mathbf{b}_1$、$2\pi\mathbf{b}_2$、$2\pi\mathbf{b}_3$ 张成。利用 (1.19)，其体积为
\begin{align*}
(2\pi)^3(\mathbf{b}_1\cdot\mathbf{b}_2\wedge\mathbf{b}_3)
&= (2\pi)^3
\frac{(\mathbf{a}_2\wedge\mathbf{a}_3)\cdot\{(\mathbf{a}_3\wedge\mathbf{a}_1)\wedge(\mathbf{a}_1\wedge\mathbf{a}_2)\}}
{(\mathbf{a}_1\cdot\mathbf{a}_2\wedge\mathbf{a}_3)^3} \\
&= (2\pi)^3
\frac{(\mathbf{a}_2\wedge\mathbf{a}_3)\cdot[\{\mathbf{a}_3\cdot(\mathbf{a}_1\wedge\mathbf{a}_2)\}\mathbf{a}_1 - \{\mathbf{a}_1\cdot(\mathbf{a}_1\wedge\mathbf{a}_2)\}\mathbf{a}_3]}
{(\mathbf{a}_1\cdot\mathbf{a}_2\wedge\mathbf{a}_3)^3} \\
&= \frac{(2\pi)^3}{(\mathbf{a}_1\cdot\mathbf{a}_2\wedge\mathbf{a}_3)} \\
&= \frac{8\pi^3}{v_c},
\tag{1.32}
\end{align*}
其中 $v_c$ 是由 $\mathbf{a}_1$、$\mathbf{a}_2$、$\mathbf{a}_3$ 张成的晶胞的体积。

这里出现因子 $8\pi^3$，是由于我们对倒格子的定义方式。更常见的约定是写成
\begin{equation*}
e^{i\boldsymbol{g}\cdot\boldsymbol{l}} \equiv \exp(2\pi i\,\mathbf{K}_g\cdot\mathbf{R}_l),
\tag{1.33}
\end{equation*}
其中 $\mathbf{R}_l$ 是格矢，而矢量 $\mathbf{K}_g$ 定义为晶胞体积为 $1/v_c$ 的另一个倒格子的矢量。这种记法的缺点是因子 $2\pi$ 总是在指数中冒出来；它也同自由电子波函数的惯用符号
\begin{equation*}
\psi = e^{i\mathbf{k}\cdot\mathbf{r}}
\tag{1.34}
\end{equation*}
不相协调。

(iv) \emph{正格子正是它自己的倒格子的倒格子。}

这已隐含在名称之中，也可以通过构造例如矢量 $(\mathbf{b}_1\wedge\mathbf{b}_2)/(\mathbf{b}_1\cdot\mathbf{b}_2\wedge\mathbf{b}_3)$ 并证明它与 $\mathbf{a}_3$ 全同来验证。考察 (1.32) 的推导过程也能看出这一点。

(v) \emph{倒格子的晶胞不必是平行六面体。}事实上，我们几乎总是在同倒格子的 Wigner--Seitz 胞（Wigner--Seitz cell）打交道。它叫做布里渊区（Brillouin zone）。

作为倒格子诸性质的一个例子，考虑一种重要且常见的结构——面心立方（face-centred cubic）格子（f.c.c. 结构）。它由四个相互穿插的简单立方格子搭成；随便看其中哪一个，都会看到它的晶胞每个面的中心各有一个原子，同时立方体的角上也分布着格点（图 10(a)）。

乍一看，这像一个带基元的格子：由 $(\mathbf{a}_x,\mathbf{a}_y,\mathbf{a}_z)$ 生成的立方晶胞里有四个原子。

但是，如果改取立方体各个面的半对角线
\begin{equation*}
\left.
\begin{aligned}
\mathbf{a}_1 &= (0, \tfrac{1}{2}a, \tfrac{1}{2}a), \qquad
\mathbf{a}_2 = (\tfrac{1}{2}a, 0, \tfrac{1}{2}a), \\
\mathbf{a}_3 &= (\tfrac{1}{2}a, \tfrac{1}{2}a, 0),
\end{aligned}
\right\}
\tag{1.35}
\end{equation*}
%FIG{10}: 图 10　面心立方格子。(a) 四个相互穿插的子晶格。(b) 作为 Bravais 格子。

那么用这些矢量的倍数的组合就能到达任何一个格点（图 10(b)）。因此，f.c.c. 格子确实是一个 Bravais 格子，$\mathbf{a}_1$、$\mathbf{a}_2$、$\mathbf{a}_3$ 就是它的生成矢。

现在我们本可以用代数方法构造出倒格子来，但利用上面指出的几何性质更为容易。例如，显然存在垂直于立方体棱边的重要点阵面。相对沿 $\mathbf{a}_x$、$\mathbf{a}_y$、$\mathbf{a}_z$ 方向的通常笛卡儿轴而言，这些面的米勒指数必为 (100)、(010)、(001) 等，即它们构成 $\{100\}$ 组。还可以看出，这些面彼此相距 $\tfrac{1}{2}a$。于是，它们对应于长度为
\begin{equation*}
|\boldsymbol{g}| = 2\pi/\tfrac{1}{2}a
\end{equation*}
沿倒格子空间中直角笛卡儿轴方向的倒格矢。我们的倒格子必须包含由这些矢量生成的整个简单立方格子。

还有一种重要的点阵面垂直于立方晶胞的对角线。这种面在三条轴上的截距相等，故其指数必属 $\{111\}$ 组。这些面
%FIG{11}: 图 11　(a) f.c.c. 格子的 {111} 面。(b) 相应的倒格子晶胞。
彼此相隔的距离等于立方体整条对角线的三分之一，即 $a/\sqrt{3}$。因此，相应的倒格矢有三个相等的分量，总长为 $2\pi\sqrt{3}/a$；它必是
\begin{equation*}
\boldsymbol{g}' = \frac{2\pi}{a}(1, 1, 1).
\tag{1.36}
\end{equation*}
在倒格子空间中，它显然沿立方晶胞对角线的方向——但长度只有对角线的一半。于是，它生成简单立方格子的各个体心。

其余的点阵面现在都可以如法研究——例如 $\{110\}$ 组面，它们垂直于基本立方体各面的对角线——不过与它们相应的倒格矢全都已经包含在由 $\{100\}$ 面和 $\{111\}$ 面生成的那些点之中。因此，f.c.c. 格子的倒格子是 b.c.c.（当然，反之亦然）。
